\songtitle{Från Stockholm till Punta Espada}

\tag*{2026}\tag{folk-rock}\tag{swedish-lyrics}\tag{thomas-entourage}\tag{erika-czajka}\tag{punta-espada}\label{sec:fran-stockholm-till-punta-espada}
% -> entourage-eco.tex
\begin{notes}
Lyrics by DeepSeek following structural instructions, narrowed themes and tone by Carlos Thompson, based on \emph{Fallstudie}, an original story by Carlos Thompson adapted to voice and tone by ChatGPT, from the \emph{Thomas' Entourage} narrative universe, narrated by Erika Czajka -- her thesis on greening a dry tropical zone with Ciudad Punta Espada as the case, Thomas's offer of help, and the journey from London and Utrecht to La Guajira. \emph{Mörkret i Punta Espada} (page \pageref{sec:morkret-i-punta-espada}) comes from the same story. Retrieved from Dvigitt's Suno page (V6-MINI, 18 November 2026).
\end{notes}
\begin{style}
Folk-rock, strummed acoustic guitar and picked acoustic figures against pulsing synth chords and warm bass; steady midtempo backbeat; tape saturation rounds the drums and thickens the chorus; Swedish female vocals, nasal and reedy across a wide range, climbing into the top notes.
\end{style}
\begin{lyrics}
\songpart{Vers 1}
Jag skrev min uppsats om grönskan i en torr och tropisk zon\\
Ciudad Punta Espada, La Guajira — en studie av vision\\
Jag visste knappt vad jag gav mig in i, bara att Thomas fanns där\\
Han sa: "Jag hjälper dig med allt jag kan" — och plötsligt var jag där

\songpart{Refräng}
Från London till Bogotá, från Utrecht till La Guajira\\
En resa genom projektet, genom människor jag fick lära känna\\
Samuel och Boris, Natalia och Victoria\\
Varje nod en lärdom, varje samtal en dörr som öppnades

\songpart{Vers 2}
London blev en särskild plats, design-teamet samlat där\\
Samuel Foreman, chefsutvecklare, visionär utan gräns\\
Boris Murphy, chefsarkitekt, teknisk och precis\\
Och Natalia, Samuels assistent — först en friktion, sedan respekt

\songpart{Refräng}
Från London till Bogotá, från Utrecht till La Guajira\\
En resa genom projektet, genom människor jag fick lära känna\\
Samuel och Boris, Natalia och Victoria\\
Varje nod en lärdom, varje samtal en dörr som öppnades

\songpart{Vers 3}
Victoria träffade jag i Barcelona, på en julfest Thomas höll\\
En vänskap som bara uppstod, där intimiteten föll\\
Så när jag var i London hade jag ingen soffa kvar\\
Jag hade en varm säng som väntade — och det var nog sårar

\songpart{Brygga}
Natalia förstod jag till slut, hennes roll var att bevaka\\
Varje galen idé jag hade måste mot visionen matcha\\
Hon var barriären Samuel behövde, för han var för konceptuell\\
Och Boris var för teknisk — Natalia höll dem samman, hon var hel

\songpart{Vers 4}
Strax efter att jag presenterat min uppsats ringde Samuel mig\\
Han insisterade: "Din vision ska få leva, inte bli liggande i en hylla"\\
Två månader efter att jag skrivit på gick Samuel bort\\
Cancern tog honom fort — och projektet stod i sorg

\songpart{Brygga}
Jag trodde Boris skulle ta över, eller någon utifrån\\
Men Natalia fick tjänsten — och det var faktiskt logiskt då\\
Hon behöver ingen storslagen vision, hon förstår helheten\\
Hon vet vad som passar in och vad som inte gör det

\songpart{Refräng}
Från London till Bogotá, från Utrecht till La Guajira\\
En resa genom projektet, genom människor jag fick lära känna\\
Samuel och Boris, Natalia och Victoria\\
Varje nod en lärdom, varje samtal en dörr som öppnades

\songpart{Outro}
Så jag skrev på, och jag stannade kvar\\
Min vision fick leva, den fick bli verklighet\\
Från en uppsats till ett projekt, från teori till mark\\
Ciudad Punta Espada — och jag är en del av det
\end{lyrics}

